% Kurzbeschreibung Strukturierter_Datenimport_LLM.py
% Maximal 30 Zeilen, Stichpunkte / kurze Sätze

\textbf{Strukturierter\_Datenimport\_LLM.py -- Freitext zu Kandidatendaten}
Nutzt ein Ollama-LLM (Standardmodell \texttt{llama2}) zur Extraktion strukturierter Kandidatendaten aus Freitext.
Definiert erlaubte Positionen und Fachbereiche (z.B. \glqq assistenzarzt\grqq{}, \glqq radiologie\grqq{}).
Beschreibt in einem detaillierten Prompt alle Zielfelder und Extraktionsregeln.
Verlangt vom LLM strikt gültiges JSON ohne zusätzliche Erklärtexte.
Erzwingt, dass nur Informationen verwendet werden, die explizit im Text vorkommen (Anti-Halluzinations-Regeln).
Definiert ein Pydantic-Modell \texttt{CandidateExtract} mit optionalen Feldern für Kerninformationen.
Enthält eine Normalisierungsfunktion, die die rohe LLM-Antwort säubert und vereinheitlicht.
Normalisiert Positionsangaben auf vordefinierte Karrierestufen (z.B. \glqq oberarzt\grqq{}).
Normalisiert Fachbereiche auf die erlaubten Fachauswahl-Begriffe.
Vereinheitlicht \texttt{regionale\_verfuegbarkeit} (z.B. Zahlen, km-Angaben, zusammengesetzte Strukturen).
Sendet den Prompt als POST-Request an die Ollama-API und liest das Feld \texttt{response}.
Schneidet aus der LLM-Antwort den JSON-Block heraus und parst ihn mit \texttt{json.loads}.
Wandelt das normalisierte Dict in eine \texttt{CandidateExtract}-Instanz um (Validierung durch Pydantic).
Löst bei ungültiger LLM-Antwort (kein JSON / Schemafehler) einen klaren Fehler aus.
Kann im Hauptblock mit Beispieltext getestet werden und liefert dann einen strukturierten Kandidatendatensatz.

